\section{Archimedean Profiles}
\label{pf:section}

In this section we prove that the profiles used in
Section~\ref{sec:certificate} are admissible in the sense of
Definition~\ref{geo:admissible} (Proposition~\ref{pf:admissible}), and we
derive bounds for the functionals $J_\R$ and $J_\C$
(Definition~\ref{geo:profiles}) that reduce to finite computations with
rational data: $J_\R$ in closed form (Lemma~\ref{pf:gaussian}), and a lower
bound for $J_\C$ (Corollary~\ref{pf:functional}) from a lower bound for the
overlap $I_\C$ (Proposition~\ref{pf:overlap}) and an upper bound for the mass
$A_\C$ (Proposition~\ref{pf:mass}). Section~\ref{sec:certificate} carries out
these computations at $\delta=\deltastar$. We use a Gaussian at the compact
coordinates, and at the pair coordinates a Student profile multiplied by a
positive polynomial.

Throughout, $0<\delta<1$ and $p=2/(1+\delta)$. We write $\psi=\Gamma'/\Gamma$
for the digamma function, $\gamma$ for Euler's constant, $(x)_n$ for the
rising factorial, $H_n=\sum_{k=1}^n1/k$ with $H_0=0$ for the harmonic numbers,
and $\mathrm{Beta}(\alpha,\alpha')$ for the beta distribution on $(0,1)$ with
density proportional to $t^{\alpha-1}(1-t)^{\alpha'-1}$. The letters $s$ and
$a$ (the parameters of Definition~\ref{pf:student}), $\psi$, $H_n$, and the
auxiliary quantities $K_n(i,k)$, $D_x(j)$, $E_{rr'}$, $E_{-,n_*}$, $G_n$,
$R_n$, $\rho$, $\Delta_1$, $\Delta_2$ introduced below are local to this
section and to Section~\ref{sec:certificate}; they are unrelated to the
complex variable $s$, the linear forms $\psi_o$ of Section~\ref{dh:section}
and the other uses of these letters elsewhere.

\begin{definition}\label{pf:student}
Let $s>0$ and $a>0$, and let $P$ be a real polynomial in two variables
that is positive on $[0,1]^2$. For $(z,w)\in\C^2$ put
\[
 t=(1+a|z|^2)^{-1},\qquad t'=(1+a|w|^2)^{-1},\qquad
 g_\C(z,w)=t^s\,t'^{\,s}P(t,t').
\]
We call $g_\C$ the \emph{Student--Bernstein profile} with parameters
$(s,a,P)$. For $P\equiv1$ it is the \emph{Student profile}
$t^st'^{\,s}$; for $s>1$, each factor $(1+a|z|^2)^{-s}$ is proportional
to the density of a bivariate Student $t$-distribution. If $P$ has
degree at most $m$ in each variable, its \emph{Bernstein coefficients}
of degree $m$ are the numbers $\beta_{ij}$ with
\[
 P(t,t')=\sum_{i,j=0}^m\beta_{ij}\,b_i^m(t)\,b_j^m(t'),\qquad
 b_i^m(t)=\binom mi t^i(1-t)^{m-i},
\]
and $b_i^m$ are the \emph{Bernstein polynomials} of degree $m$. We write
$\min\beta$ and $\max\beta$ for the smallest and largest Bernstein
coefficient.
\end{definition}

\begin{lemma}\label{pf:bernstein-bounds}
If $P$ has Bernstein coefficients $\beta_{ij}$ of degree $m$, then
$\min\beta\le P(t,t')\le\max\beta$ for $(t,t')\in[0,1]^2$.
\end{lemma}

\begin{proof}
The Bernstein polynomials $b_i^m$ are nonnegative on $[0,1]$ and sum to
one there, by the binomial theorem applied to $(t+(1-t))^m$. So $P(t,t')$ is
a convex combination of the $\beta_{ij}$, with weights
$b_i^m(t)b_j^m(t')$.
\end{proof}

From now on, $g_\C$ denotes a Student--Bernstein profile with parameters
$(s,a,P)$. Section~\ref{sec:certificate} uses
\[
 s=1.147487,\qquad a=0.0000181980309901,
\]
and the polynomial $P$ of degree $m=3$ in each variable whose Bernstein
coefficients form the symmetric matrix
\begin{equation}\label{pf:bernstein-witness}
 (\beta_{ij})=\begin{pmatrix}
 1&2.5157692640&2.8533380109&7.1926104676\\
 2.5157692640&4.1154447568&4.9424163713&8.8954745458\\
 2.8533380109&4.9424163713&5.8082616140&10.4120890697\\
 7.1926104676&8.8954745458&10.4120890697&15.3284466977
 \end{pmatrix}.
\end{equation}
All finite decimals here are exact rational numbers.

\emph{Heuristic.} The shape of $g_\C$ is adapted to the displacements
at a pair of coordinates. By \eqref{geo:norm-modulus} they have moduli
$(e^u,e^{-u})$, so as $u$ varies they run along a hyperbola. A profile
of width $a^{-1/2}$ in each coordinate stays close to its translate for
all $u$ with $e^{|u|}\lesssim a^{-1/2}$, a range of length about
$\log(1/a)$; this produces the factor $\log(1/a)$ in
Proposition~\ref{pf:overlap}, against the term $2\delta\log a$ in
\eqref{pf:collected-functional}, which comes from the exponent
$1+\delta$ on the mass $A_\C$. For Student--Bernstein profiles the
overlap $I_\C$ reduces to beta integrals and digamma values
(Lemmas~\ref{pf:laplace}--\ref{pf:pair-reduction}), which makes rigorous
bounds possible, and the polynomial $P$ couples the two radii.

\subsection{The Gaussian at the Compact Coordinates}

This subsection shows that the Gaussian with the parameter $a_\R=2p\delta$ is
optimal among Gaussians and gives $J_\R$ in closed form
(Lemma~\ref{pf:gaussian}).

\begin{lemma}\label{pf:gaussian}
Let $a_\R=2p\delta$ and $f_\R(z)=e^{-a_\R|z|^2/p}$, so that
$f_\R^p=e^{-a_\R|z|^2}$. Then
\begin{equation}\label{pf:gaussian-compact}
 J_\R=\log(p/2)+\delta\log(a_\R/\pi)-a_\R/(2p)
 =\log(p/2)+\delta\bigl(\log(a_\R/\pi)-1\bigr),
\end{equation}
and $a_\R$ maximizes $J_\R$ over the Gaussians $e^{-a'|z|^2/p}$, $a'>0$.
\end{lemma}

\begin{proof}
For $f=e^{-a'|z|^2/p}$ we have $A_\R=\pi/a'$ and, since
$|z|^2+|z+1|^2=2|z+1/2|^2+1/2$, $I_\R=(\pi p/(2a'))e^{-a'/(2p)}$. Hence
$J_\R=\log(p/2)+\delta\log(a'/\pi)-a'/(2p)$, a concave function of
$\log a'$ with its maximum at $a'=2p\delta$.
\end{proof}

By \eqref{geo:general-margin}, the margin is affine in $\theta$ with
slope $-\log\pi-2J_\R+J_\C$. This slope depends on the profiles; for
ours it is positive (Section~\ref{sec:certificate}).

\subsection{The Pair Convolution}

This subsection reduces the overlap $I_\C$ of a Student--Bernstein profile to
expectations, under beta laws, of the function $Q_{AB}$ defined below
(Lemma~\ref{pf:pair-reduction}), and expands $Q_{AB}$ in a series with
controlled remainder (Lemma~\ref{pf:hyperbola}).

\begin{lemma}\label{pf:laplace}
Let $\alpha,\alpha'>0$ with $A=\alpha+\alpha'-1>0$, and let $h\in\C$.
Then
\begin{equation}\label{pf:laplace-convolution}
 \int_{\C}(1+|z|^2)^{-\alpha}(1+|z+h|^2)^{-\alpha'}\,dz
 =\frac\pi A\,\mathbb E\bigl(1+T(1-T)|h|^2\bigr)^{-A},
 \qquad T\sim\mathrm{Beta}(\alpha,\alpha').
\end{equation}
\end{lemma}

\begin{proof}
Insert
$(1+|z|^2)^{-\alpha}=\Gamma(\alpha)^{-1}\int_0^\infty
r^{\alpha-1}e^{-r(1+|z|^2)}\,dr$ and the analogous formula with a
variable $r'$. Since
$r|z|^2+r'|z+h|^2=(r+r')|z+r'h/(r+r')|^2+rr'|h|^2/(r+r')$, the
integral of $e^{-r|z|^2-r'|z+h|^2}$ over $z\in\C$ is
$\pi(r+r')^{-1}e^{-rr'|h|^2/(r+r')}$; the factors $e^{-r}$ and
$e^{-r'}$ remain in the Laplace integrals. Substitute $r=RT$,
$r'=R(1-T)$, with Jacobian $R$, and integrate over $R$. This gives
\[
 \frac{\pi\,\Gamma(A)}{\Gamma(\alpha)\Gamma(\alpha')}
 \int_0^1T^{\alpha-1}(1-T)^{\alpha'-1}\bigl(1+T(1-T)|h|^2\bigr)^{-A}dT,
\]
which equals the right side of \eqref{pf:laplace-convolution} because
$\Gamma(\alpha+\alpha')=A\,\Gamma(A)$. All integrands are nonnegative,
so Tonelli's theorem justifies the interchanges.
\end{proof}

For $A,B>0$ and $y>0$ define
\[
 Q_{AB}(y)=\int_\R(1+ye^{2v})^{-A}(1+ye^{-2v})^{-B}\,dv .
\]

\begin{lemma}\label{pf:hyperbola}
Let $A,B>0$.
\begin{enumerate}
\item[(a)] For $y>0$,
\begin{equation}\label{pf:elementary-hyperbola-error}
 \Bigl|Q_{AB}(y)-\log(1/y)+\frac{\psi(A)+\psi(B)+2\gamma}{2}\Bigr|
 \le\frac{(A+B)y}2 .
\end{equation}
\item[(b)] For $0<y<1$, with absolutely convergent series,
\[
 Q_{AB}(y)=\sum_{n\ge0}\frac{(A)_n(B)_n}{(n!)^2}\,y^{2n}
 \Bigl[\log(1/y)+\psi(n+1)-\frac{\psi(A+n)+\psi(B+n)}2\Bigr].
\]
\item[(c)] Suppose $A,B\ge1$, $0<y<1$, $ABy^2<1$ and
$2\log(1/y)\ge H_{\lceil A\rceil-1}+H_{\lceil B\rceil-1}$. Then every
term of the series in (b) is nonnegative, and for every integer
$n_*\ge0$ its remainder $\operatorname{Rem}_{n_*}(y)$ after the terms $n\le n_*$
satisfies
\[
 0\le\operatorname{Rem}_{n_*}(y)\le\log(1/y)\frac{(ABy^2)^{n_*+1}}{1-ABy^2}.
\]
\end{enumerate}
\end{lemma}

\begin{proof}
(a) Split the integral at $v=0$. On $v\ge0$, first replace
$(1+ye^{-2v})^{-B}$ by $1$. Since
$0\le1-(1+ye^{-2v})^{-B}\le Bye^{-2v}$, this changes the half-integral
by an amount in $[0,By/2]$. The substitution $x=ye^{2v}$ gives
$\int_0^\infty(1+ye^{2v})^{-A}dv=\frac12\int_y^\infty(1+x)^{-A}dx/x$,
and
\[
 \int_y^\infty\frac{dx}{x(1+x)^A}
 =\log\frac1y+\varkappa_A+\int_0^y\bigl(1-(1+x)^{-A}\bigr)\frac{dx}x,\qquad
 \varkappa_A=\int_0^\infty\bigl((1+x)^{-A}-\mathbf 1_{x<1}\bigr)\frac{dx}x .
\]
The last integral lies in $[0,Ay]$, because $0\le1-(1+x)^{-A}\le Ax$;
in particular it tends to $0$ with $y$. On the other hand, the
substitution $t=1/(1+x)$ gives
\[
 \int_y^\infty\frac{dx}{x(1+x)^A}
 =\int_0^{1/(1+y)}\frac{t^{A-1}-1}{1-t}\,dt+\log\frac{1+y}y .
\]
Letting $y\to0$ in the two expressions yields
$\varkappa_A=\int_0^1(t^{A-1}-1)(1-t)^{-1}dt=-\psi(A)-\gamma$
\cite[Equation~(5.9.16)]{DLMF}. Thus the half $v\ge0$ equals
$\frac12(\log(1/y)-\psi(A)-\gamma)$ plus a term in $[0,Ay/2]$ minus a
term in $[0,By/2]$. The half $v\le0$ is the same with $A$ and $B$
exchanged. Adding the halves proves (a).

(b) The substitutions $x=ye^{2v}$ and then $t=x/(x+y^2)$ give
\[
 2Q_{AB}(y)=\int_0^1t^{B-1}(1-t)^{A-1}\bigl(1-(1-y^2)t\bigr)^{-A}dt .
\]
By Euler's integral \cite[Equation~(15.6.1)]{DLMF}, the right side is
$\Gamma(A)\Gamma(B)\,\mathbf F(A,B;A+B;1-y^2)$, where $\mathbf F$ is
Olver's regularized hypergeometric function. The expansion
\cite[Equation~(15.8.10)]{DLMF}, in the case where the third parameter
is the sum of the first two, gives
\[
 \mathbf F(A,B;A+B;1-y^2)=\frac1{\Gamma(A)\Gamma(B)}\sum_{n\ge0}
 \frac{(A)_n(B)_n}{(n!)^2}\,y^{2n}\bigl[2\psi(n+1)-\psi(A+n)-\psi(B+n)
 -\log y^2\bigr]
\]
for $0<y<1$. Dividing by two proves (b). The coefficients grow at most
polynomially in $n$, so the series converges absolutely.

(c) For $A\ge1$ and $n\ge0$, monotonicity of $\psi$ and the recurrence
$\psi(x+1)=\psi(x)+1/x$ give
$\psi(n+1)\le\psi(A+n)\le\psi(\lceil A\rceil+n)\le\psi(n+1)+H_{\lceil
A\rceil-1}$. Hence the bracket in (b) lies between
$\log(1/y)-\frac12(H_{\lceil A\rceil-1}+H_{\lceil B\rceil-1})\ge0$ and
$\log(1/y)$. Moreover $(A)_n/n!=\prod_{k<n}(A+k)/(k+1)\le A^n$ for
$A\ge1$, and similarly for $B$. Summing the geometric series proves
(c).
\end{proof}

\begin{definition}\label{pf:cross}
Let $d_{ij}$ be the monomial coefficients of $P$, so that
$P(t,t')=\sum_{i,j}d_{ij}t^it'^{\,j}$, and put $\eta_0=2s-1$,
\[
 A_{ik}=\eta_0+i+k,\qquad B_{jl}=\eta_0+j+l,\qquad
 v_{ijkl}=\frac{d_{ij}d_{kl}}{A_{ik}B_{jl}} .
\]
A quadruple $(i,j,k,l)$ with $d_{ij}d_{kl}\ne0$ is a \emph{cross}.
\end{definition}

\begin{lemma}\label{pf:pair-reduction}
Let $s\ge1$ and $a>0$. Then
\[
 I_\C=\frac{\pi^2}{a^2}\sum_{(i,j,k,l)}v_{ijkl}\,
 \mathbb E\,Q_{A_{ik}B_{jl}}(y),\qquad
 y=a\sqrt{T(1-T)T'(1-T')}\le a/4,
\]
where the sum runs over all crosses and, for each cross,
$T\sim\mathrm{Beta}(s+i,s+k)$ and $T'\sim\mathrm{Beta}(s+j,s+l)$ are
independent. Each expectation is finite.
\end{lemma}

\begin{proof}
Expand $g_\C(z,w)=\sum_{i,j}d_{ij}(1+a|z|^2)^{-s-i}(1+a|w|^2)^{-s-j}$ and
multiply out $g_\C(z,w)g_\C(z+e^u,w+e^{-u})$. For one cross, the scaling
$z\mapsto z/\sqrt a$ and Lemma~\ref{pf:laplace} give
\[
 \int_\C(1+a|z|^2)^{-s-i}(1+a|z+e^u|^2)^{-s-k}\,dz
 =\frac\pi{aA_{ik}}\,\mathbb E\bigl(1+aT(1-T)e^{2u}\bigr)^{-A_{ik}},
\]
and similarly in $w$ with $e^{-2u}$, $B_{jl}$ and $T'$. By Tonelli's
theorem, the $u$-integral of the product is
$\mathbb E\int_\R(1+Ye^{2u})^{-A_{ik}}(1+Y'e^{-2u})^{-B_{jl}}du$ with
$Y=aT(1-T)$ and $Y'=aT'(1-T')$, and the translation
$u\mapsto u+\frac14\log(Y'/Y)$ turns the inner integral into
$Q_{A_{ik}B_{jl}}(\sqrt{YY'})$. By
\eqref{pf:elementary-hyperbola-error}, $Q_{AB}(y)$ is at most
$\log(1/y)$ plus a constant, and $\log(1/(T(1-T)))$ is integrable under
every beta law; so each term is finite and the finite signed sum is
legitimate.
\end{proof}

\subsection{Admissibility}

This subsection proves the Fourier bound (P5) at the pair coordinates by a
contour shift (Lemma~\ref{pf:tube}), and with it the admissibility of the
profiles (Proposition~\ref{pf:admissible}). For the Fourier bound we assume that $P$ has positive Bernstein
coefficients $\beta_{ij}$ of some degree $m\ge1$.

\begin{definition}\label{pf:tube-constant}
Let $\Delta_1$ and $\Delta_2$ be the forward difference operators in
the two indices of $(\beta_{ij})$. For $0<\rho<1$, the width of the
complex tube in the proof of Lemma~\ref{pf:tube}, put
$\omega=(\rho+\rho^2)/(1-\rho^2)$,
\[
 E_{rr'}=\binom mr\binom m{r'}\max_{i,j}\bigl|\Delta_1^r\Delta_2^{r'}
 \beta_{ij}\bigr|,\qquad
 q_0=\frac1{\min\beta}\sum_{r+r'>0}E_{rr'}\,\omega^{r+r'} .
\]
We call $q_0$ the \emph{tube constant}.
\end{definition}

\begin{lemma}\label{pf:tube}
Let $s>0$, $sp>1$ and $a>0$, let $P$ have positive Bernstein coefficients
$\beta_{ij}$ of degree $m\ge1$, let $0<\rho<1$, and suppose that the tube
constant satisfies $q_0<1$. Then for all $\xi_1,\xi_2\in\C$,
\[
 \bigl|\widehat{g_\C^p}(\xi_1,\xi_2)\bigr|
 \le K_\C e^{-\varsigma_\C(|\xi_1|+|\xi_2|)}A_\C,\qquad
 K_\C=(1-\rho^2)^{-2sp}(1+q_0)^p,\qquad
 \varsigma_\C=2\pi\rho/\sqrt a .
\]
\end{lemma}

\begin{proof}
Write $x=\sqrt a\,z\in\R^2$ and consider complex vectors $x+iy$ with
$x,y\in\R^2$ and $|y|\le\rho$. Put
$\mathcal Z(x+iy)=1+(x_1+iy_1)^2+(x_2+iy_2)^2$. Then
\[
 \operatorname{Re}\mathcal Z=1+|x|^2-|y|^2\ge(1-\rho^2)(1+|x|^2),\qquad
 \bigl|\mathcal Z^{-1}-(1+|x|^2)^{-1}\bigr|
 \le\frac{2|x|\rho+\rho^2}{(1-\rho^2)(1+|x|^2)^2}\le\omega .
\]
So the complexified $t=\mathcal Z^{-1}$ lies within $\omega$ of the real
$t_0=(1+|x|^2)^{-1}\in(0,1]$, and likewise for $t'$. By the Bernstein
derivative formula,
$\partial_t^r\partial_{t'}^{r'}P/(r!\,r'!)=\binom mr\binom m{r'}
\sum_{i,j}(\Delta_1^r\Delta_2^{r'}\beta)_{ij}\,b_i^{m-r}(t)b_j^{m-r'}(t')$,
which is at most $E_{rr'}$ in absolute value on $[0,1]^2$. Taylor
expansion of the polynomial $P$ about $(t_0,t'_0)$ therefore gives
$|P(t,t')-P(t_0,t'_0)|\le q_0\min\beta\le q_0P(t_0,t'_0)$, the last step by
Lemma~\ref{pf:bernstein-bounds}. Thus
$P(t,t')$ lies in the right half-plane, and with principal branches,
\[
 \bigl|\mathcal Z(x+iy)^{-sp}\mathcal Z(x'+iy')^{-sp}P(t,t')^p\bigr|
 \le K_\C\,t_0^{sp}\,t_0'^{\,sp}P(t_0,t'_0)^p ,
\]
where $x'+iy'$ is the variable attached to $w$. The same bounds hold
for $|y|,|y'|<\rho'$ with some $\rho'>\rho$, since $q_0$ depends
continuously on $\rho$, so the left side is holomorphic in each complex
coordinate on a neighborhood of the closed tube.

The function $g_\C^p$ is invariant under rotations of $z$ and of $w$.
Rotate so that $\xi_1$ and $\xi_2$ are positive multiples of the first
basis vector, and shift the first real coordinate of $x$ by $-i\rho$,
then that of $x'$ by $-i\rho$. Each shift is justified by Cauchy's
theorem in one variable: the integrand is holomorphic in the strip and
bounded there by a constant times $(1+|x|^2)^{-sp}(1+|x'|^2)^{-sp}$,
which tends to zero at infinity and is integrable because $sp>1$. After
the shifts the Fourier factor is multiplied by
$e^{-2\pi\rho(|\xi_1|+|\xi_2|)/\sqrt a}$, and the displayed bound gives
the lemma.
\end{proof}

\begin{proposition}\label{pf:admissible}
Let $s\ge1$, $a>0$, and let $P$ have positive Bernstein coefficients
$\beta_{ij}$ of some degree $m\ge1$ in each variable. Let $0<\delta<1$ with
$a_\R=2p\delta\le1$, let $f_\R$ be the Gaussian of
Lemma~\ref{pf:gaussian}, and let $0<\rho<1$ with $q_0<1$. Then
$(f_\R,g_\C)$ is admissible (Definition~\ref{geo:admissible}) with Fourier constants
\begin{equation}\label{pf:polynomial-fourier-contract}
 M=\max\{2,\sqrt{K_\C}\},\qquad \varsigma=\min\{1,\varsigma_\C\}.
\end{equation}
\end{proposition}

\begin{proof}
(P1) The functions $f_\R^p$ and
$g_\C^p=(1+a|z|^2)^{-sp}(1+a|w|^2)^{-sp}P(t,t')^p$ are smooth, because
$P\ge\min\beta>0$ on $[0,1]^2$ (Lemma~\ref{pf:bernstein-bounds}), and they and all their derivatives are
bounded.

(P2) The integrals at the compact coordinates are computed in
Lemma~\ref{pf:gaussian}. The mass $A_\C$ is finite because $sp>1$
(Proposition~\ref{pf:mass} below), and $0<I_\C<\infty$ because
$g_\C\le\max\beta\,t^st'^{\,s}$ and Lemma~\ref{pf:pair-reduction}
applies to the Student profile $t^st'^{\,s}$.

(P3) $V_\R=a_\R|z|^2/p$, and
$V_\C=s\log(1+a|z|^2)+s\log(1+a|w|^2)-\log P(t,t')$ is at least
$-\log\max\beta$ and tends to infinity with $|z|$ or $|w|$.

(P4) The endpoint law at a compact coordinate is Gaussian. For the pair,
$-\log P$ is bounded, so it suffices to bound the second moments of
$\log(1+a|z|^2)$ and $\log(1+a|w|^2)$. Choose $0<\tau<s$; then
$\tau<2s-1$ as well, since $s\ge1$. We have
$\log^2(1+x)\le c_\tau(1+x)^\tau$ for $x\ge0$ and some constant
$c_\tau$, and $g_\C\le\max\beta\,t^st'^{\,s}$. Hence the second moment
of $\log(1+a|z|^2)$ is at most a constant times the integral of
$(1+a|z|^2)^{-(s-\tau)}(1+a|z+e^u|^2)^{-s}(1+a|w|^2)^{-s}
(1+a|w+e^{-u}|^2)^{-s}$. By the argument of
Lemma~\ref{pf:pair-reduction}, with exponents $s-\tau>0$ and $s$ and
$A=2s-\tau-1>0$, this integral is a multiple of
$\mathbb E\,Q_{A,2s-1}(y)$, which is finite by
\eqref{pf:elementary-hyperbola-error}. The same holds for $w$.

(P5) By Lemma~\ref{geo:gaussian} the Gaussian satisfies
$|\widehat{f_\R^p}(\xi)|\le2e^{-|\xi|}A_\R\le Me^{-\varsigma|\xi|}A_\R$, and by
Lemma~\ref{pf:tube} the pair profile satisfies the bound with
$K_\C e^{-\varsigma_\C(|\xi_1|+|\xi_2|)}\le
M^2e^{-\varsigma(|\xi_1|+|\xi_2|)}$.
\end{proof}

\subsection{\texorpdfstring{A Lower Bound for the Overlap $I_\C$}{A Lower Bound for the Overlap}}

This subsection bounds $I_\C$ from below by finitely many rational terms, a
logarithm and one digamma constant (Proposition~\ref{pf:overlap}).

\begin{definition}\label{pf:overlap-coefficients}
For $n\ge0$ and integers $i,k\ge0$ put
\[
 K_n(i,k)=\frac{(s+i)_n(s+k)_n}{n!\,(\eta_0+i+k+n)_{n+1}},\qquad
 D_x(j)=\sum_{h=0}^{j-1}\frac1{x+h}\quad(D_x(0)=0),
\]
\[
 r_n(i,k)=-\frac1{\eta_0}+\frac{H_n}2-\frac{D_{\eta_0}(i+k+n)}2
 +D_{\eta_0}(i+k+2n+1)-\frac{D_s(i+n)+D_s(k+n)}2,
\]
$R_n(i,k)=K_n(i,k)r_n(i,k)$, and
\[
 G_n=\sum d_{ij}d_{kl}K_n(i,k)K_n(j,l),\qquad
 \widetilde G_n=\sum d_{ij}d_{kl}\bigl[R_n(i,k)K_n(j,l)+K_n(i,k)R_n(j,l)\bigr],
\]
where the sums run over all crosses. Put
\begin{equation}\label{pf:student-constant}
 c_s=-2\psi(s)+\psi(2s)+(2s-1)^{-1}-\gamma ,\qquad y_0=a/4 .
\end{equation}
\end{definition}

For rational $s$ and $d_{ij}$, the numbers $K_n(i,k)$, $r_n(i,k)$,
$R_n(i,k)$, $G_n$ and $\widetilde G_n$ are rational.

\begin{proposition}\label{pf:overlap}
Let $s\ge1$ and $a>0$, and suppose that for every cross
\begin{equation}\label{pf:cross-threshold}
 A_{ik}B_{jl}y_0^2<1,\qquad
 2\log(4/a)\ge H_{\lceil A_{ik}\rceil-1}+H_{\lceil B_{jl}\rceil-1},
 \qquad \log(4/a)>1/2 .
\end{equation}
Then for every integer $n_*\ge0$,
\begin{equation}\label{pf:collected-overlap}
 \frac{I_\C}{\pi^2a^{-2}}\ge
 \sum_{n=0}^{n_*}a^{2n}\bigl\{G_n[\log(1/a)+c_s]+\widetilde G_n\bigr\}
 -E_{-,n_*},
\end{equation}
where the \emph{remainder term} is
\[
 E_{-,n_*}=\log(4/a)\sum_{v_{ijkl}<0}(-v_{ijkl})
 \frac{(A_{ik}B_{jl}y_0^2)^{n_*+1}}{1-A_{ik}B_{jl}y_0^2}.
\]
\end{proposition}

\begin{proof}
Start from Lemma~\ref{pf:pair-reduction}. Since $y\le y_0=a/4<1$, the
hypotheses of Lemma~\ref{pf:hyperbola}(c) hold at every value of $y$,
with $A=A_{ik}\ge1$ and $B=B_{jl}\ge1$. Split each
$Q_{A_{ik}B_{jl}}(y)$ into the terms $n\le n_*$ of
Lemma~\ref{pf:hyperbola}(b) and the remainder $\operatorname{Rem}_{n_*}(y)$.

Consider the term of index $n$ for one cross, with $\alpha=s+i$,
$\alpha'=s+k$, $A=A_{ik}=\alpha+\alpha'-1$, and similarly for $T'$. Its
factor $y^{2n}$ is $a^{2n}(T(1-T))^n(T'(1-T'))^n$, and
$\log(1/y)=\log(1/a)-\frac12\log(T(1-T))-\frac12\log(T'(1-T'))$. The
beta integral gives $\mathbb E(T(1-T))^n=(\alpha)_n(\alpha')_n/(A+1)_{2n}$,
so
\[
 \frac{(A)_n}{A\,n!}\,\mathbb E(T(1-T))^n
 =\frac{(\alpha)_n(\alpha')_n}{n!\,(A+n)_{n+1}}=K_n(i,k).
\]
Weighting by $(T(1-T))^n$ turns $\mathrm{Beta}(\alpha,\alpha')$ into
$\mathrm{Beta}(\alpha+n,\alpha'+n)$, under which the mean of
$\log(T(1-T))$ is $\psi(\alpha+n)+\psi(\alpha'+n)-2\psi(A+2n+1)$. The
recurrence $\psi(x+j)=\psi(x)+D_x(j)$ and $\psi(1)=-\gamma$ give
\[
 \begin{gathered}
 \psi(n+1)=-\gamma+H_n,\qquad
 \psi(\alpha+n)=\psi(s)+D_s(i+n),\qquad
 \psi(\alpha'+n)=\psi(s)+D_s(k+n),\\
 \psi(A+n)=\psi(\eta_0)+D_{\eta_0}(i+k+n),\qquad
 \psi(A+2n+1)=\psi(\eta_0)+D_{\eta_0}(i+k+2n+1),\\
 \tfrac12c_s=-\psi(s)+\tfrac12\psi(\eta_0)+\eta_0^{-1}-\tfrac12\gamma,
 \end{gathered}
\]
the last one because $\psi(2s)=\psi(\eta_0)+1/\eta_0$. Substituting,
\[
 \tfrac12\psi(n+1)-\tfrac12\psi(A+n)
 -\tfrac12\bigl[\psi(\alpha+n)+\psi(\alpha'+n)\bigr]+\psi(A+2n+1)
 =\tfrac12c_s+r_n(i,k).
\]
Hence the expectation of the term of index $n$, multiplied by
$v_{ijkl}$, is
\[
 a^{2n}d_{ij}d_{kl}K_n(i,k)K_n(j,l)\bigl[\log(1/a)+c_s+r_n(i,k)
 +r_n(j,l)\bigr].
\]
Summing over crosses gives the displayed main terms.

By Lemma~\ref{pf:hyperbola}(c),
$0\le\operatorname{Rem}_{n_*}(y)\le\log(1/y)(ABy^2)^{n_*+1}/(1-ABy^2)$. The right
side increases with $y$ on $(0,y_0]$, because the derivative of
$y^{2n_*+2}\log(1/y)$ is $y^{2n_*+1}((2n_*+2)\log(1/y)-1)>0$ when
$\log(1/y)>1/2$. So
$0\le\operatorname{Rem}_{n_*}(y)\le\log(4/a)(ABy_0^2)^{n_*+1}/(1-ABy_0^2)$. For
crosses with $v_{ijkl}>0$ we discard the remainder, and for crosses
with $v_{ijkl}<0$ we subtract its upper bound. This proves
\eqref{pf:collected-overlap}.
\end{proof}

With $n_*=1$ the bound keeps the exact term of order $a^2$; the terms it
discards or subtracts are of order $a^4\log(1/a)$. It is a lower bound
for $I_\C$ only.

\subsection{\texorpdfstring{An Upper Bound for the Mass $A_\C$}{An Upper Bound for the Mass}}

This subsection bounds $A_\C$ from above by subdividing $[0,1]^2$ and using
Bernstein coefficients on each rectangle (Proposition~\ref{pf:mass}).

\begin{proposition}\label{pf:mass}
Let $g_\C$ be a Student--Bernstein profile with parameters $(s,a,P)$
such that $q_\star=sp-1>0$. Then
\begin{equation}\label{pf:polynomial-mass}
 A_\C=\frac{\pi^2}{a^2q_\star^2}\,\bar A_\C,\qquad
 \bar A_\C=\mathbb E\,P(T,T')^p,
\end{equation}
where $T$ and $T'$ are independent with density $q_\star t^{q_\star-1}$ on $(0,1)$.
Partition $[0,1]^2$ into rectangles $\mathcal Q=[l,r]\times[l',r']$, and
use on $\mathcal Q$ the coordinates $x=(T-l)/(r-l)$ and
$x'=(T'-l')/(r'-l')$. Suppose that the Bernstein coefficients, of some
degree, of the restriction of $P$ to $\mathcal Q$ in these coordinates
lie in $[m_{\mathcal Q},M_{\mathcal Q}]$ with $m_{\mathcal Q}>0$, and
put
\[
 c_{\mathcal Q}=\frac{m_{\mathcal Q}+M_{\mathcal Q}}2,\qquad
 \rho_{\mathcal Q}=\frac{M_{\mathcal Q}-m_{\mathcal Q}}
 {M_{\mathcal Q}+m_{\mathcal Q}},\qquad
 P=c_{\mathcal Q}(1+W_{\mathcal Q})\ \text{on }\mathcal Q .
\]
Define
\begin{equation}\label{pf:subrectangle-moments}
 \mathrm{Mom}_i(l,r)=\int_l^rq_\star v^{q_\star-1}\Bigl(\frac{v-l}{r-l}\Bigr)^idv
 =\frac{q_\star}{(r-l)^i}\sum_{j=0}^i\binom ij(-l)^{i-j}
 \frac{r^{q_\star+j}-l^{q_\star+j}}{q_\star+j},
\end{equation}
and, writing $W_{\mathcal Q}^k=\sum_{i,j}c^{(k)}_{ij}x^ix'^{\,j}$,
\[
 \Sigma_{\mathcal Q,N}=c_{\mathcal Q}^p\sum_{k=0}^N\binom pk
 \sum_{i,j}c^{(k)}_{ij}\mathrm{Mom}_i(l,r)\mathrm{Mom}_j(l',r').
\]
Then for every $N\ge1$,
\begin{equation}\label{pf:subrectangle-mass}
 \Bigl|\bar A_\C-\sum_{\mathcal Q}\Sigma_{\mathcal Q,N}\Bigr|
 \le\sum_{\mathcal Q}c_{\mathcal Q}^p
 \frac{|\binom p{N+1}|\rho_{\mathcal Q}^{N+1}}{1-\rho_{\mathcal Q}}
 M_0(l,r)M_0(l',r').
\end{equation}
\end{proposition}

\begin{proof}
The substitution $t=(1+a|z|^2)^{-1}$ gives $dz=\pi\,dt/(at^2)$ on $\C$.
Hence
\[
 A_\C=\frac{\pi^2}{a^2}\int_0^1\!\!\int_0^1t^{sp-2}t'^{\,sp-2}
 P(t,t')^p\,dt\,dt',
\]
which is \eqref{pf:polynomial-mass} since $sp-2=q_\star-1$. Lemma~\ref{pf:bernstein-bounds}, applied on $\mathcal Q$
in the coordinates $(x,x')$, gives $m_{\mathcal Q}\le P\le M_{\mathcal Q}$, that
is, $|W_{\mathcal Q}|\le\rho_{\mathcal Q}<1$ on $\mathcal Q$. For $1<p<2$ the binomial coefficients satisfy
$|\binom p{k+1}|/|\binom pk|=|p-k|/(k+1)<1$ for $k\ge1$, so the binomial
series of $(1+W)^p$ converges uniformly on $|W|\le\rho_{\mathcal Q}$,
and its tail after index $N\ge1$ is at most
$|\binom p{N+1}|\rho_{\mathcal Q}^{N+1}/(1-\rho_{\mathcal Q})$.
Integrate against the density on $\mathcal Q$, whose integral over
$\mathcal Q$ is $M_0(l,r)M_0(l',r')$. The moment formula follows by expanding
$(v-l)^i$; the convention $0^{q_\star}=0$ is valid since $q_\star>0$.
\end{proof}

\begin{remark}\label{pf:symmetry}
If $P$ is symmetric, that is, $P(t,t')=P(t',t)$, and both axes are
partitioned alike, then the reflection $(T,T')\mapsto(T',T)$ preserves the
density of Proposition~\ref{pf:mass} and maps each rectangle $\mathcal Q$ of
the partition to a rectangle $\mathcal Q'$ of the partition, with the same
Bernstein coefficients up to transposition, hence the same $c_{\mathcal Q}$
and $\rho_{\mathcal Q}$. So $\mathcal Q$ and $\mathcal Q'$ contribute equally
to the sum and to the error bound in \eqref{pf:subrectangle-mass}, and
off-diagonal rectangles may be evaluated once and counted twice, diagonal
ones once. Subsection~\ref{cert:computations} describes how the quantities in
Proposition~\ref{pf:mass} are evaluated.
\end{remark}

\subsection{\texorpdfstring{A Lower Bound for the Functional $J_\C$}{A Lower Bound for the Functional}}

This subsection combines the two previous ones into the lower bound for $J_\C$
used in Definition~\ref{cert:lower-bounds} (Corollary~\ref{pf:functional}),
and encloses the digamma values that it needs (Lemma~\ref{pf:digamma}).

\begin{corollary}\label{pf:functional}
Under the hypotheses of Propositions~\ref{pf:overlap} and~\ref{pf:mass},
let $Z_*$ be the right side of \eqref{pf:collected-overlap} for some
$n_*$, and suppose $Z_*>0$. If $\bar A_\C^*\ge\bar A_\C$, then
\begin{equation}\label{pf:collected-functional}
 J_\C\ge\log Z_*+2(1+\delta)\log q_\star-2\delta\log\pi
 +2\delta\log a-(1+\delta)\log\bar A_\C^* .
\end{equation}
\end{corollary}

\begin{proof}
By Proposition~\ref{pf:overlap}, $I_\C\ge\pi^2a^{-2}Z_*$, and by
\eqref{pf:polynomial-mass}, $A_\C\le\pi^2a^{-2}q_\star^{-2}\bar A_\C^*$.
Substitute into $J_\C=\log I_\C-(1+\delta)\log A_\C$.
\end{proof}

Interval evaluation of the right side of \eqref{pf:collected-functional}
gives a rigorous lower bound for $J_\C$; the upper endpoint of such an
enclosure is not an upper bound for $J_\C$. For rational $s$, all
digamma values in Proposition~\ref{pf:overlap} reduce to rational
corrections and to $c_s=2/\eta_0+\psi(\eta_0)-2\psi(s)+\psi(1)$, since
$\psi(2s)=\psi(\eta_0)+1/\eta_0$ and $\psi(1)=-\gamma$. The digamma
function at rational arguments is enclosed as follows.

\begin{lemma}\label{pf:digamma}
For $x>0$ and integers $L\ge0$ and $k_*\ge1$, put $z=x+L$. Then
\[
 \Bigl|\psi(x)-\Bigl(\log z-\frac1{2z}
 -\sum_{k=1}^{k_*-1}\frac{\mathrm B_{2k}}{2kz^{2k}}
 -\sum_{j=0}^{L-1}\frac1{x+j}\Bigr)\Bigr|
 \le\frac{|\mathrm B_{2k_*}|}{2k_*z^{2k_*}},
\]
where $\mathrm B_{2k}$ are the Bernoulli numbers.
\end{lemma}

\begin{proof}
By Binet's formula \cite[Equation~(5.9.15)]{DLMF},
\[
 \psi(z)=\log z-\frac1{2z}
 -2\int_0^\infty\frac{t\,dt}{(t^2+z^2)(e^{2\pi t}-1)} .
\]
Write
\[
 \frac1{t^2+z^2}=\sum_{k=1}^{k_*-1}(-1)^{k-1}\frac{t^{2k-2}}{z^{2k}}
 +(-1)^{k_*-1}\frac{t^{2k_*-2}}{z^{2k_*-2}(t^2+z^2)},
\]
and use $\int_0^\infty t^{2k-1}(e^{2\pi t}-1)^{-1}dt=|\mathrm B_{2k}|/(4k)$ and
$(-1)^{k-1}|\mathrm B_{2k}|=\mathrm B_{2k}$. The remaining integral has a fixed sign and
is at most $|\mathrm B_{2k_*}|/(2k_*z^{2k_*})$ in absolute value, because
$(t^2+z^2)^{-1}\le z^{-2}$. Finally apply the recurrence
$\psi(x+1)=\psi(x)+1/x$ $L$ times.
\end{proof}
